% --- Archon physics formalization source begin ---
% archon:physics
% archon:covers PhyXMiniProblems/problem_phyx_mini_0202.lean
% archon:source-report reports/phyx_mini/problem_phyx_mini_0202.source.json

\chapter{Physics problem phyx\_mini\_0202}
\label{ch:PhyXMiniProblems_problem_phyx_mini_0202}

\paragraph{Problem source.}
\#\# Physical scenario

Agent Arlene devised the following method of measuring the muzzle velocity of a rifle (figure). She fires a bullet into a 4.148\textbackslash{},\textbackslash{}mathrm\{kg\} wooden block resting on a smooth surface, and attached to a spring of spring constant \textbackslash{}( k = 162.7\textbackslash{},\textbackslash{}mathrm\{N/m\} \textbackslash{}). The bullet, whose mass is 7.870\textbackslash{},\textbackslash{}mathrm\{g\}, remains embedded in the wooden block. She measures the maximum distance that the block compresses the spring to be 9.460\textbackslash{},\textbackslash{}mathrm\{cm\}.

\#\# Question

What is the speed \textbackslash{}( v \textbackslash{}) of the bullet?

\#\# Answer choices

- A: A: 302.6m/s

- B: B: 332.6m/s

- C: C: 322.6m/s

- D: D: 312.6m/s

\#\# Figure caption (auxiliary; use the image as primary evidence)

The image depicts two diagrams illustrating a physical scenario involving a spring and colliding masses.

1. **Top Diagram:**
   - A bullet or small object with mass \textbackslash{}( m \textbackslash{}) is depicted on the left, moving towards the right.
   - It has a velocity vector \textbackslash{}( \textbackslash{}vec\{v\} \textbackslash{}) pointing towards a larger block.
   - The larger block, labeled with mass \textbackslash{}( M \textbackslash{}), is positioned horizontally on a surface.
   - Attached to the block is a spring with a spring constant \textbackslash{}( k \textbackslash{}), fixed to a wall on the right.

2. **Bottom Diagram:**
   - After the collision, the bullet is embedded in the block. The combined mass is labeled \textbackslash{}( M + m \textbackslash{}).
   - This configuration is on the same horizontal surface.
   - The spring is compressed.
   - The compression distance is labeled as 9.460 cm.

The setup likely demonstrates a collision and subsequent compression of the spring, illustrating principles of momentum conservation and energy transfer.

\#\# Dataset metadata

Momentum and Collisions; Multi-Formula Reasoning, Physical Model Grounding Reasoning

\paragraph{Recorded answer/context.}
D: D: 312.6m/s

\paragraph{Figure/image path.}
/root/proposal\_for\_physic/science-mango/phyx\_mini\_run/phyx\_data/test\_image/202.png

\paragraph{Formalization target.}
create a compiling Lean file with sorry bodies at `PhyXMiniProblems/problem\_phyx\_mini\_0202.lean`. The Lean declarations must preserve the physical quantities, dimensions or dimensional roles, figure labels, governing-law hypotheses, and final relation expressed by this problem.
Use Mathlib/Physlib names found through LeanExplore where available. If a physics API is missing, introduce faithful local abstractions rather than scalar placeholder aliases.

\begin{theorem}[Physics formalization target]
\label{thm:physics:phyx_mini_0202:target}
\lean{PhyXMiniProblems.ProblemPhyXMini0202.problem_phyx_mini_0202}
\uses{def:physics:phyx-mini-0202:phyxminiproblems-problemphyxmini0202-ballisticspringsetup, def:physics:phyx-mini-0202:phyxminiproblems-problemphyxmini0202-matchesproblemreadouts, def:physics:phyx-mini-0202:phyxminiproblems-problemphyxmini0202-matchesscenarioandsuppliedfigure, def:physics:phyx-mini-0202:phyxminiproblems-problemphyxmini0202-hasphysicalballisticspringparameters, def:physics:phyx-mini-0202:phyxminiproblems-problemphyxmini0202-satisfiescollisionmomentumconservation, def:physics:phyx-mini-0202:phyxminiproblems-problemphyxmini0202-satisfiesenergyconservationduringcompression, def:physics:phyx-mini-0202:phyxminiproblems-problemphyxmini0202-roundstonearesttenthmeterpersecond, def:physics:phyx-mini-0202:phyxminiproblems-problemphyxmini0202-matchesanswerchoice}
The assigned autoformalize agent should translate this physics problem into Lean declarations in the covered file, with theorem and lemma proof bodies written as `by sorry`.
\end{theorem}
\begin{proof}
This is an autoformalization task, not a proof task. Produce statements that can later be proved by the physics prover without weakening the physical model.
\end{proof}

% --- Archon declaration topology begin ---
\section{Declaration topology}
\begin{definition}[Mass Quantity]
  \label{def:physics:phyx-mini-0202:phyxminiproblems-problemphyxmini0202-massquantity}
  \lean{PhyXMiniProblems.ProblemPhyXMini0202.MassQuantity}
  A nonnegative physical mass, independent of its readout unit
\end{definition}
\begin{proof}
  This is the declared model definition of Mass Quantity.
\end{proof}
\begin{definition}[Length Quantity]
  \label{def:physics:phyx-mini-0202:phyxminiproblems-problemphyxmini0202-lengthquantity}
  \lean{PhyXMiniProblems.ProblemPhyXMini0202.LengthQuantity}
  A nonnegative physical length, independent of its readout unit
\end{definition}
\begin{proof}
  This is the declared model definition of Length Quantity.
\end{proof}
\begin{definition}[Speed Quantity]
  \label{def:physics:phyx-mini-0202:phyxminiproblems-problemphyxmini0202-speedquantity}
  \lean{PhyXMiniProblems.ProblemPhyXMini0202.SpeedQuantity}
  A nonnegative speed, with dimension length per time
\end{definition}
\begin{proof}
  This is the declared model definition of Speed Quantity.
\end{proof}
\begin{definition}[Spring Stiffness Quantity]
  \label{def:physics:phyx-mini-0202:phyxminiproblems-problemphyxmini0202-springstiffnessquantity}
  \lean{PhyXMiniProblems.ProblemPhyXMini0202.SpringStiffnessQuantity}
  Spring stiffness, with dimension mass per time squared (N/m)
\end{definition}
\begin{proof}
  This is the declared model definition of Spring Stiffness Quantity.
\end{proof}
\begin{definition}[mass Readout]
  \label{def:physics:phyx-mini-0202:phyxminiproblems-problemphyxmini0202-massreadout}
  \lean{PhyXMiniProblems.ProblemPhyXMini0202.massReadout}
  \uses{def:physics:phyx-mini-0202:phyxminiproblems-problemphyxmini0202-massquantity}
  Read a physical mass in a selected mass unit
\end{definition}
\begin{proof}
  This is the declared model definition of mass Readout.
\end{proof}
\begin{definition}[length Readout]
  \label{def:physics:phyx-mini-0202:phyxminiproblems-problemphyxmini0202-lengthreadout}
  \lean{PhyXMiniProblems.ProblemPhyXMini0202.lengthReadout}
  \uses{def:physics:phyx-mini-0202:phyxminiproblems-problemphyxmini0202-lengthquantity}
  Read a physical length in a selected length unit
\end{definition}
\begin{proof}
  This is the declared model definition of length Readout.
\end{proof}
\begin{definition}[speed Readout]
  \label{def:physics:phyx-mini-0202:phyxminiproblems-problemphyxmini0202-speedreadout}
  \lean{PhyXMiniProblems.ProblemPhyXMini0202.speedReadout}
  \uses{def:physics:phyx-mini-0202:phyxminiproblems-problemphyxmini0202-speedquantity}
  Read a speed in the coherent length-per-time unit selected below
\end{definition}
\begin{proof}
  This is the declared model definition of speed Readout.
\end{proof}
\begin{definition}[stiffness Readout]
  \label{def:physics:phyx-mini-0202:phyxminiproblems-problemphyxmini0202-stiffnessreadout}
  \lean{PhyXMiniProblems.ProblemPhyXMini0202.stiffnessReadout}
  \uses{def:physics:phyx-mini-0202:phyxminiproblems-problemphyxmini0202-springstiffnessquantity}
  Read spring stiffness in the coherent mass-per-time-squared unit
\end{definition}
\begin{proof}
  This is the declared model definition of stiffness Readout.
\end{proof}
\begin{definition}[mass In Kilograms]
  \label{def:physics:phyx-mini-0202:phyxminiproblems-problemphyxmini0202-massinkilograms}
  \lean{PhyXMiniProblems.ProblemPhyXMini0202.massInKilograms}
  \uses{def:physics:phyx-mini-0202:phyxminiproblems-problemphyxmini0202-massquantity, def:physics:phyx-mini-0202:phyxminiproblems-problemphyxmini0202-massreadout}
  Kilogram readout used for the two masses and the governing laws
\end{definition}
\begin{proof}
  This is the declared model definition of mass In Kilograms.
\end{proof}
\begin{definition}[mass In Grams]
  \label{def:physics:phyx-mini-0202:phyxminiproblems-problemphyxmini0202-massingrams}
  \lean{PhyXMiniProblems.ProblemPhyXMini0202.massInGrams}
  \uses{def:physics:phyx-mini-0202:phyxminiproblems-problemphyxmini0202-massquantity, def:physics:phyx-mini-0202:phyxminiproblems-problemphyxmini0202-massreadout}
  Gram readout used for the bullet mass stated in the problem
\end{definition}
\begin{proof}
  This is the declared model definition of mass In Grams.
\end{proof}
\begin{definition}[length In Meters]
  \label{def:physics:phyx-mini-0202:phyxminiproblems-problemphyxmini0202-lengthinmeters}
  \lean{PhyXMiniProblems.ProblemPhyXMini0202.lengthInMeters}
  \uses{def:physics:phyx-mini-0202:phyxminiproblems-problemphyxmini0202-lengthquantity, def:physics:phyx-mini-0202:phyxminiproblems-problemphyxmini0202-lengthreadout}
  Meter readout used for spring displacements
\end{definition}
\begin{proof}
  This is the declared model definition of length In Meters.
\end{proof}
\begin{definition}[length In Centimeters]
  \label{def:physics:phyx-mini-0202:phyxminiproblems-problemphyxmini0202-lengthincentimeters}
  \lean{PhyXMiniProblems.ProblemPhyXMini0202.lengthInCentimeters}
  \uses{def:physics:phyx-mini-0202:phyxminiproblems-problemphyxmini0202-lengthquantity, def:physics:phyx-mini-0202:phyxminiproblems-problemphyxmini0202-lengthreadout}
  Centimeter readout used for the measured maximum compression
\end{definition}
\begin{proof}
  This is the declared model definition of length In Centimeters.
\end{proof}
\begin{definition}[speed In Meters Per Second]
  \label{def:physics:phyx-mini-0202:phyxminiproblems-problemphyxmini0202-speedinmeterspersecond}
  \lean{PhyXMiniProblems.ProblemPhyXMini0202.speedInMetersPerSecond}
  \uses{def:physics:phyx-mini-0202:phyxminiproblems-problemphyxmini0202-speedquantity, def:physics:phyx-mini-0202:phyxminiproblems-problemphyxmini0202-speedreadout}
  Speed readout in meters per SI second
\end{definition}
\begin{proof}
  This is the declared model definition of speed In Meters Per Second.
\end{proof}
\begin{definition}[stiffness In Newtons Per Meter]
  \label{def:physics:phyx-mini-0202:phyxminiproblems-problemphyxmini0202-stiffnessinnewtonspermeter}
  \lean{PhyXMiniProblems.ProblemPhyXMini0202.stiffnessInNewtonsPerMeter}
  \uses{def:physics:phyx-mini-0202:phyxminiproblems-problemphyxmini0202-springstiffnessquantity, def:physics:phyx-mini-0202:phyxminiproblems-problemphyxmini0202-stiffnessreadout}
  Spring-stiffness readout in kilograms per second squared, i.e. newtons per meter
\end{definition}
\begin{proof}
  This is the declared model definition of stiffness In Newtons Per Meter.
\end{proof}
\begin{definition}[Figure Panel]
  \label{def:physics:phyx-mini-0202:phyxminiproblems-problemphyxmini0202-figurepanel}
  \lean{PhyXMiniProblems.ProblemPhyXMini0202.FigurePanel}
  The two moments depicted by the upper and lower panels of the image
\end{definition}
\begin{proof}
  This is the declared model definition of Figure Panel.
\end{proof}
\begin{definition}[Figure Object]
  \label{def:physics:phyx-mini-0202:phyxminiproblems-problemphyxmini0202-figureobject}
  \lean{PhyXMiniProblems.ProblemPhyXMini0202.FigureObject}
  Objects distinguished in the supplied two-panel figure
\end{definition}
\begin{proof}
  This is the declared model definition of Figure Object.
\end{proof}
\begin{definition}[Figure Label]
  \label{def:physics:phyx-mini-0202:phyxminiproblems-problemphyxmini0202-figurelabel}
  \lean{PhyXMiniProblems.ProblemPhyXMini0202.FigureLabel}
  Mathematical labels printed in the supplied figure
\end{definition}
\begin{proof}
  This is the declared model definition of Figure Label.
\end{proof}
\begin{definition}[Horizontal Direction]
  \label{def:physics:phyx-mini-0202:phyxminiproblems-problemphyxmini0202-horizontaldirection}
  \lean{PhyXMiniProblems.ProblemPhyXMini0202.HorizontalDirection}
  Horizontal direction of the bullet's velocity arrow in the upper panel
\end{definition}
\begin{proof}
  This is the declared model definition of Horizontal Direction.
\end{proof}
\begin{definition}[Collision Outcome]
  \label{def:physics:phyx-mini-0202:phyxminiproblems-problemphyxmini0202-collisionoutcome}
  \lean{PhyXMiniProblems.ProblemPhyXMini0202.CollisionOutcome}
  Qualitative collision outcome represented in the lower panel
\end{definition}
\begin{proof}
  This is the declared model definition of Collision Outcome.
\end{proof}
\begin{definition}[Surface Condition]
  \label{def:physics:phyx-mini-0202:phyxminiproblems-problemphyxmini0202-surfacecondition}
  \lean{PhyXMiniProblems.ProblemPhyXMini0202.SurfaceCondition}
  Surface idealization stated in the prose
\end{definition}
\begin{proof}
  This is the declared model definition of Surface Condition.
\end{proof}
\begin{definition}[Ballistic Spring Figure]
  \label{def:physics:phyx-mini-0202:phyxminiproblems-problemphyxmini0202-ballisticspringfigure}
  \lean{PhyXMiniProblems.ProblemPhyXMini0202.BallisticSpringFigure}
  \uses{def:physics:phyx-mini-0202:phyxminiproblems-problemphyxmini0202-figurepanel, def:physics:phyx-mini-0202:phyxminiproblems-problemphyxmini0202-figureobject, def:physics:phyx-mini-0202:phyxminiproblems-problemphyxmini0202-figurelabel, def:physics:phyx-mini-0202:phyxminiproblems-problemphyxmini0202-horizontaldirection}
  Qualitative evidence transcribed from the primary image. The numerical compression datum remains a readout of the physical displacement in the setup; showsLabel only records that its label appears in the figure
\end{definition}
\begin{proof}
  This is the declared model definition of Ballistic Spring Figure.
\end{proof}
\begin{definition}[Ballistic Spring Setup]
  \label{def:physics:phyx-mini-0202:phyxminiproblems-problemphyxmini0202-ballisticspringsetup}
  \lean{PhyXMiniProblems.ProblemPhyXMini0202.BallisticSpringSetup}
  \uses{def:physics:phyx-mini-0202:phyxminiproblems-problemphyxmini0202-massquantity, def:physics:phyx-mini-0202:phyxminiproblems-problemphyxmini0202-lengthquantity, def:physics:phyx-mini-0202:phyxminiproblems-problemphyxmini0202-speedquantity, def:physics:phyx-mini-0202:phyxminiproblems-problemphyxmini0202-springstiffnessquantity, def:physics:phyx-mini-0202:phyxminiproblems-problemphyxmini0202-collisionoutcome, def:physics:phyx-mini-0202:phyxminiproblems-problemphyxmini0202-surfacecondition, def:physics:phyx-mini-0202:phyxminiproblems-problemphyxmini0202-ballisticspringfigure}
  All independent quantities used in the experiment. In particular, initialBulletSpeed is an unconstrained physical speed until the data and governing laws are supplied; no field assigns it an answer-choice value
\end{definition}
\begin{proof}
  This is the declared model definition of Ballistic Spring Setup.
\end{proof}
\begin{definition}[Matches Problem Readouts]
  \label{def:physics:phyx-mini-0202:phyxminiproblems-problemphyxmini0202-matchesproblemreadouts}
  \lean{PhyXMiniProblems.ProblemPhyXMini0202.MatchesProblemReadouts}
  \uses{def:physics:phyx-mini-0202:phyxminiproblems-problemphyxmini0202-massinkilograms, def:physics:phyx-mini-0202:phyxminiproblems-problemphyxmini0202-massingrams, def:physics:phyx-mini-0202:phyxminiproblems-problemphyxmini0202-lengthinmeters, def:physics:phyx-mini-0202:phyxminiproblems-problemphyxmini0202-lengthincentimeters, def:physics:phyx-mini-0202:phyxminiproblems-problemphyxmini0202-speedinmeterspersecond, def:physics:phyx-mini-0202:phyxminiproblems-problemphyxmini0202-stiffnessinnewtonspermeter, def:physics:phyx-mini-0202:phyxminiproblems-problemphyxmini0202-ballisticspringsetup}
  Numerical readouts stated in the problem, together with the initial and final boundary conditions implicit in "resting" and "maximum compression". The requested muzzle speed does not occur here
\end{definition}
\begin{proof}
  This is the declared model definition of Matches Problem Readouts.
\end{proof}
\begin{definition}[Matches Scenario And Supplied Figure]
  \label{def:physics:phyx-mini-0202:phyxminiproblems-problemphyxmini0202-matchesscenarioandsuppliedfigure}
  \lean{PhyXMiniProblems.ProblemPhyXMini0202.MatchesScenarioAndSuppliedFigure}
  \uses{def:physics:phyx-mini-0202:phyxminiproblems-problemphyxmini0202-ballisticspringsetup}
  Qualitative scenario and primary-image facts. The upper panel places the right-moving bullet to the left of the block and the block to the left of the wall-mounted spring. The lower panel shows the embedded composite and the compression annotation
\end{definition}
\begin{proof}
  This is the declared model definition of Matches Scenario And Supplied Figure.
\end{proof}
\begin{definition}[Has Physical Ballistic Spring Parameters]
  \label{def:physics:phyx-mini-0202:phyxminiproblems-problemphyxmini0202-hasphysicalballisticspringparameters}
  \lean{PhyXMiniProblems.ProblemPhyXMini0202.HasPhysicalBallisticSpringParameters}
  \uses{def:physics:phyx-mini-0202:phyxminiproblems-problemphyxmini0202-massinkilograms, def:physics:phyx-mini-0202:phyxminiproblems-problemphyxmini0202-lengthinmeters, def:physics:phyx-mini-0202:phyxminiproblems-problemphyxmini0202-speedinmeterspersecond, def:physics:phyx-mini-0202:phyxminiproblems-problemphyxmini0202-stiffnessinnewtonspermeter, def:physics:phyx-mini-0202:phyxminiproblems-problemphyxmini0202-ballisticspringsetup}
  Positivity and nondegeneracy conditions for the physical experiment
\end{definition}
\begin{proof}
  This is the declared model definition of Has Physical Ballistic Spring Parameters.
\end{proof}
\begin{definition}[Satisfies Collision Momentum Conservation]
  \label{def:physics:phyx-mini-0202:phyxminiproblems-problemphyxmini0202-satisfiescollisionmomentumconservation}
  \lean{PhyXMiniProblems.ProblemPhyXMini0202.SatisfiesCollisionMomentumConservation}
  \uses{def:physics:phyx-mini-0202:phyxminiproblems-problemphyxmini0202-massreadout, def:physics:phyx-mini-0202:phyxminiproblems-problemphyxmini0202-speedreadout, def:physics:phyx-mini-0202:phyxminiproblems-problemphyxmini0202-ballisticspringsetup}
  One-dimensional momentum conservation during the short perfectly inelastic collision. It is stated in every coherent selection of mass, length, and time units and retains the initially stationary block term explicitly
\end{definition}
\begin{proof}
  This is the declared model definition of Satisfies Collision Momentum Conservation.
\end{proof}
\begin{definition}[Satisfies Energy Conservation During Compression]
  \label{def:physics:phyx-mini-0202:phyxminiproblems-problemphyxmini0202-satisfiesenergyconservationduringcompression}
  \lean{PhyXMiniProblems.ProblemPhyXMini0202.SatisfiesEnergyConservationDuringCompression}
  \uses{def:physics:phyx-mini-0202:phyxminiproblems-problemphyxmini0202-massreadout, def:physics:phyx-mini-0202:phyxminiproblems-problemphyxmini0202-lengthreadout, def:physics:phyx-mini-0202:phyxminiproblems-problemphyxmini0202-speedreadout, def:physics:phyx-mini-0202:phyxminiproblems-problemphyxmini0202-stiffnessreadout, def:physics:phyx-mini-0202:phyxminiproblems-problemphyxmini0202-ballisticspringsetup}
  Mechanical-energy conservation from immediately after impact until maximum compression. Multiplying the usual kinetic-plus-spring-potential equation by two gives the displayed relation. Collision energy is deliberately not conserved: this law applies only after the embedded composite has formed
\end{definition}
\begin{proof}
  This is the declared model definition of Satisfies Energy Conservation During Compression.
\end{proof}
\begin{lemma}[initial Bullet Speed squared relation]
  \label{lem:physics:phyx-mini-0202:phyxminiproblems-problemphyxmini0202-initialbulletspeed-squared-relation}
  \lean{PhyXMiniProblems.ProblemPhyXMini0202.initialBulletSpeed_squared_relation}
  \uses{def:physics:phyx-mini-0202:phyxminiproblems-problemphyxmini0202-massinkilograms, def:physics:phyx-mini-0202:phyxminiproblems-problemphyxmini0202-lengthinmeters, def:physics:phyx-mini-0202:phyxminiproblems-problemphyxmini0202-speedinmeterspersecond, def:physics:phyx-mini-0202:phyxminiproblems-problemphyxmini0202-stiffnessinnewtonspermeter, def:physics:phyx-mini-0202:phyxminiproblems-problemphyxmini0202-ballisticspringsetup, def:physics:phyx-mini-0202:phyxminiproblems-problemphyxmini0202-matchesproblemreadouts, def:physics:phyx-mini-0202:phyxminiproblems-problemphyxmini0202-hasphysicalballisticspringparameters, def:physics:phyx-mini-0202:phyxminiproblems-problemphyxmini0202-satisfiescollisionmomentumconservation, def:physics:phyx-mini-0202:phyxminiproblems-problemphyxmini0202-satisfiesenergyconservationduringcompression}
  Eliminating the intermediate composite speed yields the exact squared-speed relation. This is a derived conclusion, not a premise of the main theorem
\end{lemma}
\begin{proof}
  The conclusion follows from the governing relations and figure readouts named in the statement of initial Bullet Speed squared relation.
\end{proof}
\begin{lemma}[initial Bullet Speed exact readout]
  \label{lem:physics:phyx-mini-0202:phyxminiproblems-problemphyxmini0202-initialbulletspeed-exact-readout}
  \lean{PhyXMiniProblems.ProblemPhyXMini0202.initialBulletSpeed_exact_readout}
  \uses{def:physics:phyx-mini-0202:phyxminiproblems-problemphyxmini0202-speedinmeterspersecond, def:physics:phyx-mini-0202:phyxminiproblems-problemphyxmini0202-ballisticspringsetup, def:physics:phyx-mini-0202:phyxminiproblems-problemphyxmini0202-matchesproblemreadouts, def:physics:phyx-mini-0202:phyxminiproblems-problemphyxmini0202-hasphysicalballisticspringparameters, def:physics:phyx-mini-0202:phyxminiproblems-problemphyxmini0202-satisfiescollisionmomentumconservation, def:physics:phyx-mini-0202:phyxminiproblems-problemphyxmini0202-satisfiesenergyconservationduringcompression}
  The positive solution of the squared relation, with all source measurements converted to SI units. The decimal answer is obtained only after rounding this exact radical expression
\end{lemma}
\begin{proof}
  The conclusion follows from the governing relations and figure readouts named in the statement of initial Bullet Speed exact readout.
\end{proof}
\begin{definition}[Answer Choice]
  \label{def:physics:phyx-mini-0202:phyxminiproblems-problemphyxmini0202-answerchoice}
  \lean{PhyXMiniProblems.ProblemPhyXMini0202.AnswerChoice}
  Labels attached to the four speeds printed in the problem
\end{definition}
\begin{proof}
  This is the declared model definition of Answer Choice.
\end{proof}
\begin{definition}[displayed Answer Speed In Meters Per Second]
  \label{def:physics:phyx-mini-0202:phyxminiproblems-problemphyxmini0202-displayedanswerspeedinmeterspersecond}
  \lean{PhyXMiniProblems.ProblemPhyXMini0202.displayedAnswerSpeedInMetersPerSecond}
  \uses{def:physics:phyx-mini-0202:phyxminiproblems-problemphyxmini0202-answerchoice}
  Displayed answer speeds, read in meters per second
\end{definition}
\begin{proof}
  This is the declared model definition of displayed Answer Speed In Meters Per Second.
\end{proof}
\begin{definition}[recorded Dataset Answer]
  \label{def:physics:phyx-mini-0202:phyxminiproblems-problemphyxmini0202-recordeddatasetanswer}
  \lean{PhyXMiniProblems.ProblemPhyXMini0202.recordedDatasetAnswer}
  \uses{def:physics:phyx-mini-0202:phyxminiproblems-problemphyxmini0202-answerchoice}
  The dataset's recorded answer label; this metadata is not a theorem premise
\end{definition}
\begin{proof}
  This is the declared model definition of recorded Dataset Answer.
\end{proof}
\begin{definition}[Rounds To Nearest Tenth Meter Per Second]
  \label{def:physics:phyx-mini-0202:phyxminiproblems-problemphyxmini0202-roundstonearesttenthmeterpersecond}
  \lean{PhyXMiniProblems.ProblemPhyXMini0202.RoundsToNearestTenthMeterPerSecond}
  \uses{def:physics:phyx-mini-0202:phyxminiproblems-problemphyxmini0202-speedquantity, def:physics:phyx-mini-0202:phyxminiproblems-problemphyxmini0202-speedinmeterspersecond}
  Agreement with a displayed speed after rounding to the nearest 0.1 m/s
\end{definition}
\begin{proof}
  This is the declared model definition of Rounds To Nearest Tenth Meter Per Second.
\end{proof}
\begin{definition}[Matches Answer Choice]
  \label{def:physics:phyx-mini-0202:phyxminiproblems-problemphyxmini0202-matchesanswerchoice}
  \lean{PhyXMiniProblems.ProblemPhyXMini0202.MatchesAnswerChoice}
  \uses{def:physics:phyx-mini-0202:phyxminiproblems-problemphyxmini0202-ballisticspringsetup, def:physics:phyx-mini-0202:phyxminiproblems-problemphyxmini0202-answerchoice, def:physics:phyx-mini-0202:phyxminiproblems-problemphyxmini0202-displayedanswerspeedinmeterspersecond, def:physics:phyx-mini-0202:phyxminiproblems-problemphyxmini0202-roundstonearesttenthmeterpersecond}
  The physical muzzle speed rounds to the number printed beside a choice
\end{definition}
\begin{proof}
  This is the declared model definition of Matches Answer Choice.
\end{proof}
% --- Archon declaration topology end ---
% --- Archon physics formalization source end ---
